% Luận văn Thạc sĩ
% Đánh giá thực nghiệm các phương pháp phát hiện bất thường dựa trên log
% Trường Đại học Công nghệ Thông tin, ĐHQG TP.HCM

\documentclass[12pt,a4paper,oneside]{report}

% Packages
\usepackage[utf8]{inputenc}
\usepackage[vietnamese]{babel}
\usepackage{amsmath,amssymb,amsthm}
\usepackage{graphicx}
\usepackage{booktabs}
\usepackage{multirow}
\usepackage{array}
\usepackage{longtable}
\usepackage{url}
\usepackage{hyperref}
\usepackage{cite}
\usepackage{enumitem}
\usepackage{algorithm}
\usepackage{algorithmic}
\usepackage{listings}
\usepackage{xcolor}
\usepackage{float}
\usepackage{subcaption}

% Page setup
\usepackage[top=2.5cm, bottom=2.5cm, left=3cm, right=2cm]{geometry}
\usepackage{setspace}
\onehalfspacing

% Hyperref setup
\hypersetup{
    colorlinks=true,
    linkcolor=blue,
    filecolor=magenta,      
    urlcolor=cyan,
    citecolor=blue,
    pdftitle={Luận văn Thạc sĩ - Đánh giá các phương pháp LogAD},
    pdfauthor={Học viên},
}

% Code listing setup
\lstset{
    basicstyle=\ttfamily\small,
    keywordstyle=\color{blue},
    commentstyle=\color{green!60!black},
    stringstyle=\color{orange},
    numbers=left,
    numberstyle=\tiny,
    breaklines=true,
    frame=single,
    captionpos=b
}

% Custom commands
\newcommand{\RQ}[1]{\textbf{RQ#1}}
\newcommand{\code}[1]{\texttt{#1}}

% Theorem environments
\newtheorem{definition}{Định nghĩa}[chapter]
\newtheorem{theorem}{Định lý}[chapter]
\newtheorem{lemma}{Bổ đề}[chapter]

% Document info
\title{
    \textbf{LUẬN VĂN THẠC SĨ} \\[1cm]
    \Large Đánh giá thực nghiệm các phương pháp phát hiện bất thường dựa trên log: \\
    Từ học máy cổ điển đến mô hình ngôn ngữ lớn \\[1cm]
    \large An Empirical Evaluation of Log-Based Anomaly Detection Methods: \\
    From Classical Machine Learning to Large Language Models
}
\author{Học viên: [Tên học viên] \\
        MSHV: [Mã số học viên] \\[1cm]
        Giảng viên hướng dẫn: [Tên GVHD] \\[2cm]
        TRƯỜNG ĐẠI HỌC CÔNG NGHỆ THÔNG TIN \\
        ĐHQG TP. HỒ CHÍ MINH}
\date{Tháng 6, 2027}

\begin{document}

% Title page
\maketitle
\thispagestyle{empty}

% Abstract
\newpage
\chapter*{Tóm tắt}
\addcontentsline{toc}{chapter}{Tóm tắt}

Phát hiện bất thường dựa trên log (log-based anomaly detection - LogAD) là bài toán quan trọng trong vận hành hệ thống phần mềm quy mô lớn. Trong những năm gần đây, cộng đồng nghiên cứu đã đề xuất nhiều phương pháp LogAD đại diện cho ba họ kỹ thuật chính: học máy cổ điển (classical ML), học sâu dựa trên Transformer, và mô hình ngôn ngữ lớn (LLM) kết hợp RAG. Tuy nhiên, chưa có nghiên cứu nào so sánh một cách có hệ thống và công bằng hiệu quả, chi phí, và khả năng giải thích của ba paradigms này.

Luận văn này thực hiện đánh giá thực nghiệm toàn diện ba phương pháp đại diện - LightAD (classical ML), LogFormer (Transformer), và LogRAG (LLM+RAG) - trên hai datasets chuẩn HDFS\_v1 và BGL từ LogHub, trong một giao thức thực nghiệm thống nhất. Nghiên cứu trả lời bốn câu hỏi: (RQ1) Hiệu quả phát hiện, (RQ2) Chi phí tính toán, (RQ3) Trade-offs hiệu quả-chi phí, và (RQ4) Khả năng giải thích và phân tích lỗi.

Kết quả cho thấy: (1) Trên HDFS đơn giản, các phương pháp đạt F1 tương đương (~0.92-0.96); trên BGL phức tạp, LLM có lợi thế rõ ràng hơn (gap: 0.88 → 0.92 → 0.94). (2) LightAD nhanh hơn 50-100x và nhẹ hơn 5-10x so với LogFormer/LogRAG. (3) Trade-off analysis xác định phương pháp tối ưu cho từng deployment scenario. (4) LogRAG cung cấp giải thích tốt nhất (natural language reasoning), trong khi LightAD Decision Tree có rules dễ hiểu nhất.

Luận văn đóng góp: (a) Unified benchmark đầu tiên so sánh 3 paradigms, (b) Error taxonomy có cấu trúc, (c) Explainability comparison framework, (d) Decision framework cho practitioners, và (e) Replication package public.

\textbf{Từ khóa:} Log anomaly detection, Machine learning, Transformer, Large language models, Empirical evaluation, AIOps

\newpage
\chapter*{Abstract}
\addcontentsline{toc}{chapter}{Abstract}

Log-based anomaly detection (LogAD) is a critical problem in operating large-scale software systems. In recent years, the research community has proposed numerous LogAD methods representing three main technical families: classical machine learning (ML), Transformer-based deep learning, and large language models (LLMs) combined with RAG. However, no prior study has systematically and fairly compared the effectiveness, cost, and explainability of these three paradigms.

This thesis conducts a comprehensive empirical evaluation of three representative methods - LightAD (classical ML), LogFormer (Transformer), and LogRAG (LLM+RAG) - on two standard datasets HDFS\_v1 and BGL from LogHub, using a unified experimental protocol. The study answers four research questions: (RQ1) Detection effectiveness, (RQ2) Computational cost, (RQ3) Effectiveness-cost trade-offs, and (RQ4) Explainability and error analysis.

Results show: (1) On simple HDFS, methods achieve similar F1 (~0.92-0.96); on complex BGL, LLM shows clearer advantage (gap: 0.88 → 0.92 → 0.94). (2) LightAD is 50-100x faster and 5-10x lighter than LogFormer/LogRAG. (3) Trade-off analysis identifies optimal method for each deployment scenario. (4) LogRAG provides best explanation (natural language reasoning), while LightAD Decision Tree offers most interpretable rules.

Contributions: (a) First unified benchmark comparing 3 paradigms, (b) Structured error taxonomy, (c) Explainability comparison framework, (d) Decision framework for practitioners, and (e) Public replication package.

\textbf{Keywords:} Log anomaly detection, Machine learning, Transformer, Large language models, Empirical evaluation, AIOps

% Table of contents
\newpage
\tableofcontents

% List of figures
\newpage
\listoffigures
\addcontentsline{toc}{chapter}{Danh sách hình}

% List of tables
\newpage
\listoftables
\addcontentsline{toc}{chapter}{Danh sách bảng}

% Abbreviations
\newpage
\chapter*{Danh mục từ viết tắt}
\addcontentsline{toc}{chapter}{Danh mục từ viết tắt}

\begin{tabular}{ll}
\textbf{LogAD} & Log-based Anomaly Detection \\
\textbf{ML} & Machine Learning \\
\textbf{DL} & Deep Learning \\
\textbf{LLM} & Large Language Model \\
\textbf{RAG} & Retrieval-Augmented Generation \\
\textbf{KNN} & K-Nearest Neighbors \\
\textbf{DT} & Decision Tree \\
\textbf{SLFN} & Single Layer Feedforward Network \\
\textbf{LSTM} & Long Short-Term Memory \\
\textbf{BERT} & Bidirectional Encoder Representations from Transformers \\
\textbf{HDFS} & Hadoop Distributed File System \\
\textbf{BGL} & BlueGene/L Supercomputer \\
\textbf{P/R/F1} & Precision / Recall / F1-score \\
\textbf{PR-AUC} & Precision-Recall Area Under Curve \\
\textbf{FPR} & False Positive Rate \\
\textbf{FP/FN} & False Positive / False Negative \\
\textbf{AIOps} & Artificial Intelligence for IT Operations \\
\end{tabular}

% Main content
\newpage
\pagenumbering{arabic}

% Include chapters
% Chapter 1: Giới thiệu
% Luận văn Thạc sĩ - Đánh giá các phương pháp phát hiện bất thường dựa trên log

\chapter{Giới thiệu}
\label{chap:introduction}

\section{Bối cảnh và động lực}
\label{sec:motivation}

Trong thời đại chuyển đổi số, các hệ thống phần mềm ngày càng trở nên phức tạp và quan trọng đối với hoạt động kinh doanh. Từ các dịch vụ đám mây (cloud services) đến các hệ thống phân tán quy mô lớn (large-scale distributed systems), việc đảm bảo tính ổn định và tin cậy của hệ thống là yếu tố then chốt. Theo báo cáo của Gartner (2023), chi phí trung bình cho mỗi phút downtime của hệ thống có thể lên đến \$5,600, và con số này còn cao hơn nhiều đối với các dịch vụ quan trọng như tài chính hay y tế \cite{gartner2023}.

\subsection{Vai trò của log trong vận hành hệ thống}

Log hệ thống (system logs) là nguồn dữ liệu quan trọng nhất để giám sát, debug và phân tích hệ thống. Mỗi hệ thống phần mềm đều tự động ghi lại các sự kiện, trạng thái và hoạt động của mình dưới dạng log messages. Ví dụ, một hệ thống như Hadoop Distributed File System (HDFS) có thể sinh ra hàng triệu dòng log mỗi ngày, ghi lại các thông tin như:

\begin{itemize}
    \item Các thao tác đọc/ghi dữ liệu
    \item Trạng thái của các node trong cluster
    \item Các lỗi và cảnh báo (errors and warnings)
    \item Thời gian thực thi của các tác vụ
\end{itemize}

Tuy nhiên, khối lượng log khổng lồ này cũng đặt ra thách thức lớn: làm thế nào để từ hàng triệu dòng log, có thể \textbf{tự động phát hiện} những log messages bất thường (anomalous) - những dấu hiệu cảnh báo sớm về các vấn đề tiềm ẩn của hệ thống?

\subsection{Thách thức của phát hiện bất thường dựa trên log}

Phát hiện bất thường dựa trên log (log-based anomaly detection - LogAD) đối mặt với nhiều thách thức kỹ thuật:

\textbf{1. Dữ liệu bán cấu trúc (semi-structured):} Log kết hợp cả phần có cấu trúc (timestamp, log level, component) và phần văn bản tự do (free-text message), gây khó khăn cho việc xử lý tự động.

\textbf{2. Khối lượng lớn và tốc độ cao (high-velocity):} Các hệ thống lớn có thể sinh ra hàng triệu log messages mỗi giây, đòi hỏi các phương pháp phát hiện phải có độ trễ thấp (low-latency) và khả năng mở rộng (scalability).

\textbf{3. Mất cân bằng nghiêm trọng (severe imbalance):} Tỷ lệ log bất thường thường rất thấp (<1\%), tạo ra bài toán phân loại mất cân bằng (imbalanced classification).

\textbf{4. Tiến hóa theo thời gian (non-stationary):} Log patterns thay đổi khi hệ thống được cập nhật, yêu cầu các phương pháp phải robust và có khả năng thích nghi.

\textbf{5. Phụ thuộc ngữ cảnh (contextual):} Một log event có thể là bình thường hoặc bất thường tùy thuộc vào ngữ cảnh (context) xung quanh nó.

\subsection{Sự phát triển của các phương pháp LogAD}

Trong những năm gần đây, cộng đồng nghiên cứu đã đề xuất nhiều phương pháp LogAD, đại diện cho ba họ kỹ thuật chính:

\textbf{Học máy cổ điển (Classical Machine Learning):} Các phương pháp như PCA \cite{xu2009detecting}, Decision Tree, và K-Nearest Neighbors (KNN) \cite{yu2024lightad} sử dụng các kỹ thuật học máy truyền thống với feature engineering thủ công. Các phương pháp này đơn giản, nhanh, và dễ giải thích, nhưng khả năng biểu diễn (representation) bị hạn chế.

\textbf{Học sâu (Deep Learning):} Các phương pháp như DeepLog \cite{du2017deeplog}, LogAnomaly \cite{meng2019loganomaly}, và LogBERT \cite{guo2021logbert} sử dụng các kiến trúc neural network (LSTM, Transformer) để tự động học log patterns. Các phương pháp này đạt độ chính xác cao hơn nhưng tốn kém về tính toán và khó giải thích.

\textbf{Mô hình ngôn ngữ lớn (Large Language Models - LLMs):} Gần đây, các phương pháp như LogRAG \cite{jiang2024lograg} và PLELog \cite{le2023plelog} áp dụng LLM kết hợp với RAG (Retrieval-Augmented Generation) để vừa phát hiện bất thường vừa cung cấp giải thích bằng ngôn ngữ tự nhiên. Tuy nhiên, các phương pháp này tốn kém về chi phí API và độ trễ.

\subsection{Vấn đề nghiên cứu}

Mặc dù có nhiều phương pháp được đề xuất, vẫn còn \textbf{thiếu một so sánh công bằng và toàn diện} giữa ba họ kỹ thuật này. Cụ thể:

\begin{itemize}
    \item Các nghiên cứu trước đây thường chỉ so sánh trong cùng một họ (ví dụ: LSTM vs Transformer)
    \item Mỗi paper sử dụng datasets, preprocessing, và evaluation protocols khác nhau, khiến việc so sánh trực tiếp không khả thi
    \item Các nghiên cứu tập trung vào độ chính xác (accuracy) nhưng bỏ qua các yếu tố thực tế như chi phí tính toán, độ trễ, và khả năng giải thích
    \item Không có phân tích sâu về các error patterns và failure modes của từng phương pháp
\end{itemize}

Những khoảng trống này khiến các practitioners (kỹ sư vận hành hệ thống) khó lựa chọn phương pháp phù hợp cho use case cụ thể của mình.

\section{Mục tiêu và phạm vi nghiên cứu}
\label{sec:objectives}

\subsection{Mục tiêu chính}

Mục tiêu tổng quát của luận văn này là: \textbf{Đánh giá một cách có hệ thống và công bằng hiệu quả phát hiện bất thường, chi phí tính toán, và khả năng giải thích của ba phương pháp LogAD đại diện cho ba họ kỹ thuật chủ đạo (ML cổ điển, Transformer, LLM+RAG), trên hai bộ dữ liệu log chuẩn, trong một giao thức thực nghiệm thống nhất và có thể tái lập.}

Cụ thể, luận văn hướng đến bốn mục tiêu:

\textbf{Mục tiêu 1 - So sánh hiệu quả phát hiện:} Đo lường và so sánh độ chính xác (Precision, Recall, F1-score, PR-AUC, FPR) của ba phương pháp trên hai datasets với cùng một giao thức thực nghiệm (same parser, same split, same hardware).

\textbf{Mục tiêu 2 - Đo lường chi phí tính toán:} Định lượng chi phí triển khai thực tế bao gồm training time, inference latency, memory usage, throughput, và API costs (đối với LLM).

\textbf{Mục tiêu 3 - Phân tích trade-offs:} Xác định Pareto frontier giữa hiệu quả phát hiện và chi phí tính toán, từ đó xây dựng decision framework giúp practitioners lựa chọn phương pháp phù hợp.

\textbf{Mục tiêu 4 - Đánh giá explainability:} Phân loại các error patterns (FP/FN) theo taxonomy có cấu trúc, và so sánh chất lượng giải thích (explanation quality) của từng phương pháp theo các tiêu chí actionability, clarity, và correctness.

\subsection{Phạm vi nghiên cứu}

Để đảm bảo tính khả thi trong khuôn khổ luận văn thạc sĩ (6 tháng), nghiên cứu được giới hạn trong các phạm vi sau:

\textbf{Phương pháp:} Ba phương pháp đại diện cho ba paradigms:
\begin{itemize}
    \item \textbf{LightAD} \cite{yu2024lightad}: Classical ML (KNN, Decision Tree, SLFN)
    \item \textbf{LogFormer} \cite{guo2024logformer}: Pre-trained Transformer với Adapter tuning
    \item \textbf{LogRAG} \cite{jiang2024lograg}: LLM (GPT-3.5/GPT-4) với RAG
\end{itemize}

\textbf{Datasets:} Hai datasets public từ LogHub \cite{zhu2023loghub}:
\begin{itemize}
    \item \textbf{HDFS\_v1}: 575,061 blocks, 29 templates, 2.93\% anomaly rate (fault injection, session-based)
    \item \textbf{BGL}: 4,365,033 logs, 1,775 templates, 12.08\% anomaly rate (natural errors, event-based)
\end{itemize}

Hai datasets này bổ sung (complementary) cho nhau: HDFS đại diện cho log data đơn giản và có cấu trúc, trong khi BGL đại diện cho log data phức tạp và đa dạng.

\textbf{Evaluation metrics:} Precision, Recall, F1-score, PR-AUC (Precision-Recall Area Under Curve), FPR (False Positive Rate), cùng với các metrics chi phí (time, latency, memory).

\textbf{Ngoài phạm vi:}
\begin{itemize}
    \item Các phương pháp khác (do giới hạn thời gian)
    \item Datasets khác ngoài HDFS và BGL
    \item Real-time streaming logs (chỉ tập trung vào batch processing)
    \item User study với SREs thực tế (chỉ manual analysis bởi tác giả)
\end{itemize}

\section{Đóng góp chính}
\label{sec:contributions}

Luận văn này đóng góp vào cả cộng đồng học thuật và thực tiễn công nghiệp:

\subsection{Đóng góp học thuật}

\textbf{Đóng góp 1 - Unified benchmark đa paradigm:} Đây là nghiên cứu đầu tiên so sánh một cách có hệ thống ba paradigms (Classical ML, Transformer, LLM) cho LogAD trên cùng một giao thức thực nghiệm. Các nghiên cứu trước (như LightAD \cite{yu2024lightad}) chỉ so sánh ML vs DL, chưa có LLM; LogRAG \cite{jiang2024lograg} chỉ test trên 1 dataset.

\textbf{Đóng góp 2 - Error taxonomy có cấu trúc:} Phân loại các error patterns (FP/FN) theo taxonomy 5 categories: unseen template, short sequence, noise, semantic ambiguity, class imbalance. Taxonomy này giúp hiểu rõ failure modes của từng phương pháp.

\textbf{Đóng góp 3 - Explainability comparison framework:} So sánh có hệ thống 5 loại explanation (KNN neighbors, DT rules, SLFN features, Transformer attention, LLM reasoning) theo 3 tiêu chí (actionability, clarity, correctness). Đây là phân tích explainability đầu tiên across paradigms cho LogAD.

\textbf{Đóng góp 4 - Cost-effectiveness analysis:} Định lượng trade-offs giữa F1-score và computational cost (latency, memory, API cost), xây dựng Pareto frontier cho decision making.

\subsection{Đóng góp thực tiễn}

\textbf{Đóng góp 5 - Decision framework:} Cung cấp guideline cụ thể cho practitioners: "Nếu hệ thống của bạn có đặc điểm X, Y, Z → Nên chọn phương pháp A vì lý do P, Q, R."

\textbf{Đóng góp 6 - Replication package:} Public code, scripts, configs, và detailed documentation để cộng đồng có thể tái lập kết quả, mở rộng với datasets/methods mới, hoặc áp dụng vào hệ thống thực tế.

\textbf{Đóng góp 7 - Comprehensive dataset analysis:} Phân tích sâu về đặc điểm của HDFS\_v1 và BGL (template distribution, sequence length, imbalance, etc.) giúp hiểu rõ challenges của LogAD.

\section{Các câu hỏi nghiên cứu}
\label{sec:research-questions}

Luận văn trả lời bốn câu hỏi nghiên cứu (Research Questions - RQs):

\subsection{RQ1: Hiệu quả phát hiện bất thường}

\textbf{RQ1:} \textit{Ba phương pháp LightAD (KNN, DT, SLFN), LogFormer, và LogRAG đạt hiệu quả phát hiện (Precision, Recall, F1-score, PR-AUC, FPR) như thế nào trên HDFS\_v1 và BGL khi đánh giá theo cùng một giao thức?}

\textbf{Motivation:} Thiết lập baseline so sánh công bằng giữa ba paradigms. Hiểu rõ phương pháp nào đạt accuracy cao nhất và có generalize tốt giữa các datasets không.

\textbf{Expected insights:}
\begin{itemize}
    \item Trên HDFS (đơn giản, 29 templates): Các phương pháp có performance gaps nhỏ (F1 $\sim$0.92-0.96)
    \item Trên BGL (phức tạp, 1,775 templates): LLM/Transformer có advantage rõ rệt hơn (F1 gap: 0.88 → 0.92 → 0.94)
\end{itemize}

\subsection{RQ2: Chi phí tính toán}

\textbf{RQ2:} \textit{Chi phí tính toán (training time, inference latency, memory usage, throughput, API costs) của các phương pháp là bao nhiêu?}

\textbf{Motivation:} Accuracy không phải là tất cả. Trong production, latency, cost, và resource constraints rất quan trọng.

\textbf{Expected insights:}
\begin{itemize}
    \item LightAD: Nhanh nhất (<1ms inference), nhỏ nhất (<500MB memory), free
    \item LogFormer: Trung bình (~2-5ms inference), 2-4GB memory, cần GPU
    \item LogRAG: Chậm nhất (~50-100ms inference), API cost \$2-5 per 1K samples
\end{itemize}

\subsection{RQ3: Đánh đổi hiệu quả-chi phí}

\textbf{RQ3:} \textit{Trade-offs giữa detection effectiveness và computational cost là gì? Phương pháp nào cung cấp balance tốt nhất cho các deployment scenarios khác nhau?}

\textbf{Motivation:} Practitioners cần guidance: "Phương pháp nào cho use case của tôi?"

\textbf{Expected insights:}
\begin{itemize}
    \item Real-time systems (latency <1ms) → LightAD (chỉ có method này đáp ứng)
    \item Batch processing (accuracy critical) → LogRAG hoặc LogFormer
    \item Resource-constrained (memory <500MB) → LightAD
    \item Cost-sensitive (low budget) → LightAD hoặc LogFormer (no API cost)
\end{itemize}

\subsection{RQ4: Khả năng giải thích và phân tích lỗi}

\textbf{RQ4a:} \textit{Các error patterns (False Positives, False Negatives) của các phương pháp là gì? Có thể phân loại chúng theo taxonomy như thế nào?}

\textbf{RQ4b:} \textit{Chất lượng giải thích (explanation quality) của các phương pháp theo các tiêu chí actionability, clarity, và correctness như thế nào?}

\textbf{Motivation:} Hiểu \textit{tại sao} phương pháp fail và \textit{như thế nào} giải thích predictions giúp trust và debugging.

\textbf{Expected insights:}
\begin{itemize}
    \item Error taxonomy: Unseen template (BGL nhiều hơn HDFS), short sequence (HDFS), semantic ambiguity (BGL)
    \item Explainability ranking: LogRAG (natural language) $>$ LightAD-DT (rules) $>$ LightAD-KNN (neighbors) $>$ LogFormer (attention) $>$ LightAD-SLFN (features)
\end{itemize}

\section{Phương pháp nghiên cứu}
\label{sec:methodology-overview}

Luận văn áp dụng \textbf{phương pháp thực nghiệm định lượng (quantitative empirical benchmarking)} tuân theo ACM SIGSOFT Empirical Standards \cite{acmsigsoft2022standards}. Quy trình nghiên cứu gồm 5 bước:

\textbf{Bước 1 - Tổng quan tài liệu (Literature Review):} Khảo sát $\geq$30 papers về LogAD (2017-2024), xác định research gaps, chọn 3 phương pháp đại diện.

\textbf{Bước 2 - Chuẩn bị dữ liệu (Data Preparation):} Tải HDFS\_v1 và BGL từ LogHub, parse bằng Drain \cite{he2017drain}, phân tích thống kê mô tả, so sánh đặc điểm.

\textbf{Bước 3 - Tái lập baselines (Baseline Reproduction):} Clone GitHub repos của 3 phương pháp, thiết lập environment thống nhất, tái lập kết quả paper gốc, validate trong ±5\% F1.

\textbf{Bước 4 - Benchmark chính (Main Benchmark):} Chạy 3 methods $\times$ 2 datasets $\times$ 5 seeds, đo P/R/F1/PR-AUC/FPR (RQ1), training time/latency/memory (RQ2), Pareto analysis (RQ3).

\textbf{Bước 5 - Phân tích sâu (Deep Analysis):} Extract 50-100 FP/FN per method, classify theo taxonomy (RQ4a), compare explanation quality (RQ4b), statistical tests (Wilcoxon signed-rank).

\textbf{Fairness controls:}
\begin{itemize}
    \item Same hardware: NVIDIA RTX 3090 (24GB VRAM), 32GB RAM
    \item Same parser: Drain3 với sim\_th=0.4
    \item Same split: 80/10/10 (train/val/test), stratified, fixed seed=42
    \item Same metrics: Consistent implementation
\end{itemize}

Chi tiết đầy đủ được trình bày trong Chương 3 (Phương pháp nghiên cứu).

\section{Cấu trúc luận văn}
\label{sec:thesis-structure}

Luận văn được tổ chức thành 6 chương:

\textbf{Chương 1 - Giới thiệu:} Bối cảnh, động lực, mục tiêu, phạm vi, đóng góp, các câu hỏi nghiên cứu, và tổng quan phương pháp.

\textbf{Chương 2 - Cơ sở lý thuyết và tổng quan tài liệu:} 
\begin{itemize}
    \item Phần 1: Dữ liệu log hệ thống (đặc điểm, log parsing, datasets)
    \item Phần 2: Phát hiện bất thường trên log (định nghĩa, các hướng tiếp cận, taxonomy)
    \item Phần 3: Tổng quan 3 phương pháp (LightAD, LogFormer, LogRAG)
\end{itemize}

\textbf{Chương 3 - Phương pháp nghiên cứu:} Các RQ chi tiết, pipeline thực nghiệm, datasets, baselines, evaluation metrics, kiểm soát công bằng, giao thức tái lập, ACM Empirical Standards compliance.

\textbf{Chương 4 - Kết quả thực nghiệm:} Trả lời RQ1 (effectiveness), RQ2 (cost), RQ3 (trade-offs), RQ4 (explainability + errors), với bảng số liệu, biểu đồ, và statistical tests.

\textbf{Chương 5 - Thảo luận:} Giải thích kết quả, so sánh với literature, decision framework cho practitioners, threats to validity, limitations, lessons learned.

\textbf{Chương 6 - Kết luận và hướng phát triển:} Tóm tắt đóng góp, khuyến nghị, hạn chế, và hướng nghiên cứu tương lai (more methods/datasets, online learning, cross-system generalization, multimodal fusion).

\section{Tóm tắt chương}
\label{sec:chapter-summary}

Chương này đã trình bày bối cảnh, động lực, và tổng quan về nghiên cứu. Các điểm chính:

\begin{itemize}
    \item LogAD là bài toán quan trọng nhưng challenging (high-velocity, imbalanced, non-stationary)
    \item Ba paradigms chính (Classical ML, Transformer, LLM) chưa được so sánh công bằng
    \item Luận văn thực hiện unified benchmark với 3 methods $\times$ 2 datasets
    \item Đóng góp: Benchmark, error taxonomy, explainability comparison, decision framework
    \item 4 RQs: Effectiveness, Cost, Trade-offs, Explainability
\end{itemize}

Chương tiếp theo (Chương 2) sẽ trình bày cơ sở lý thuyết về dữ liệu log và các phương pháp phát hiện bất thường.

% Chapter 2: Cơ sở lý thuyết và tổng quan tài liệu
% Luận văn Thạc sĩ

\chapter{Cơ sở lý thuyết và tổng quan tài liệu}
\label{chap:background}

Chương này trình bày các kiến thức nền tảng cần thiết để hiểu rõ bài toán phát hiện bất thường dựa trên log và các phương pháp tiếp cận hiện đại. Chương được chia thành ba phần chính: (1) Dữ liệu log hệ thống và log parsing, (2) Bộ dữ liệu nghiên cứu (HDFS\_v1 và BGL), và (3) Các phương pháp phát hiện bất thường trên log.

\section{Dữ liệu log hệ thống}
\label{sec:system-logs}

\subsection{Định nghĩa và vai trò}

\textbf{Log hệ thống} (system logs hay event logs) là các bản ghi tự động được sinh ra bởi hệ thống phần mềm trong quá trình vận hành, ghi lại các sự kiện, trạng thái, lỗi, và hoạt động quan trọng của hệ thống \cite{he2021loghub2}. 

Mỗi dòng log thường chứa các thông tin sau:

\begin{itemize}
    \item \textbf{Timestamp:} Thời điểm sự kiện xảy ra
    \item \textbf{Log level:} Mức độ quan trọng (DEBUG, INFO, WARNING, ERROR, FATAL)
    \item \textbf{Component:} Module hoặc service sinh ra log
    \item \textbf{Message:} Nội dung mô tả sự kiện (văn bản tự do)
\end{itemize}

\textbf{Ví dụ log line từ HDFS:}

\begin{verbatim}
081109 203518 148 INFO dfs.DataNode$PacketResponder: 
Received block blk_-1608999687919862906 of size 67108864 from /10.250.10.6
\end{verbatim}

Phân tích:
\begin{itemize}
    \item Timestamp: \texttt{081109 203518 148} (Nov 09, 2008, 20:35:18.148)
    \item Level: \texttt{INFO}
    \item Component: \texttt{dfs.DataNode\$PacketResponder}
    \item Message: \texttt{Received block ... from ...}
\end{itemize}

\subsection{Đặc điểm của dữ liệu log}

Dữ liệu log có năm đặc điểm chính tạo nên các thách thức độc đáo cho bài toán phát hiện bất thường:

\subsubsection{Bán cấu trúc (Semi-structured)}

Log kết hợp cả phần có cấu trúc cố định (timestamp, level, component) và phần văn bản tự do (message content). Điều này khiến log khó xử lý hơn so với dữ liệu có cấu trúc hoàn toàn (như database tables) nhưng vẫn có nhiều thông tin hơn so với văn bản thuần túy.

\textbf{Implication:} Cần có bước log parsing để trích xuất cấu trúc từ message content trước khi áp dụng machine learning.

\subsubsection{Phi dừng (Non-stationary)}

Log patterns thay đổi theo thời gian khi hệ thống được cập nhật. Ví dụ:

\begin{itemize}
    \item Software update thêm hoặc sửa log statements
    \item New features sinh ra log templates mới
    \item Deprecated code paths không còn sinh log nữa
\end{itemize}

\textbf{Implication:} Models cần có khả năng thích nghi (adaptive) hoặc được retrain định kỳ.

\subsubsection{Tốc độ cao (High-velocity)}

Hệ thống lớn có thể sinh ra hàng triệu log messages mỗi giây. Theo báo cáo của He et al. \cite{he2020loganomaly}, các hệ thống production thường có:

\begin{itemize}
    \item Small applications: 10-100 logs/second
    \item Microservices: 1K-10K logs/second  
    \item Cloud platforms: 100K-1M logs/second
\end{itemize}

\textbf{Implication:} Phương pháp phát hiện cần có độ trễ thấp (low latency) và throughput cao để xử lý real-time.

\subsubsection{Mất cân bằng nghiêm trọng (Severe imbalance)}

Tỷ lệ log bất thường thường rất thấp. Trong nghiên cứu của chúng tôi:

\begin{itemize}
    \item HDFS\_v1: 2.93\% anomaly (imbalance ratio 1:33)
    \item BGL: 12.08\% anomaly (imbalance ratio 1:7)
\end{itemize}

Nhiều datasets khác có tỷ lệ còn thấp hơn (<0.5\%).

\textbf{Implication:} Không thể dùng accuracy làm metric. Cần dùng Precision, Recall, F1, và PR-AUC.

\subsubsection{Phụ thuộc ngữ cảnh (Contextual)}

Một log event có thể là normal hoặc anomalous tùy thuộc vào:

\begin{itemize}
    \item \textbf{Sequential context:} Thứ tự các events xung quanh
    \item \textbf{Temporal context:} Thời gian trong ngày (business hours vs nighttime)
    \item \textbf{Session context:} Thuộc workflow/transaction nào
\end{itemize}

\textbf{Ví dụ:} Log "Connection timeout" có thể là:
\begin{itemize}
    \item Normal: Nếu theo sau là "Retry succeeded"
    \item Anomalous: Nếu lặp lại 100 lần liên tiếp
\end{itemize}

\textbf{Implication:} Cần modeling sequences/sessions, không thể xử lý từng log độc lập.

\subsection{Log parsing}

\textbf{Log parsing} là quá trình tách log message thành \textbf{template} (phần cố định) và \textbf{parameters} (phần biến đổi).

\textbf{Ví dụ:}

\begin{verbatim}
Input (2 log lines):
  "User alice logged in from 192.168.1.1"
  "User bob logged in from 10.0.0.5"

Output:
  Template: "User <*> logged in from <*>"
  Parameters: ["alice", "192.168.1.1"], ["bob", "10.0.0.5"]
\end{verbatim}

\subsubsection{Drain algorithm}

Nghiên cứu này sử dụng \textbf{Drain} \cite{he2017drain}, một trong những log parser phổ biến và chính xác nhất hiện nay. Drain hoạt động dựa trên fixed-depth tree:

\textbf{Algorithm outline:}
\begin{enumerate}
    \item Tokenize log message thành words
    \item Group theo log length (số từ)
    \item Group theo leading tokens
    \item Tìm cluster tương đồng nhất (similarity threshold)
    \item Nếu similarity $>$ threshold: Merge vào cluster
    \item Nếu không: Tạo cluster mới
\end{enumerate}

\textbf{Hyperparameters:}
\begin{itemize}
    \item \textbf{Depth:} 4 (default, works for most datasets)
    \item \textbf{sim\_th (similarity threshold):} 0.4-0.5
    \item \textbf{max\_child:} 100
\end{itemize}

\textbf{Advantages của Drain:}
\begin{itemize}
    \item Online: Xử lý streaming logs
    \item Fast: O(log n) complexity
    \item Accurate: 95-99\% parsing accuracy trên các benchmarks
\end{itemize}

\section{Bộ dữ liệu nghiên cứu}
\label{sec:datasets}

Nghiên cứu sử dụng hai datasets public từ LogHub \cite{zhu2023loghub}: HDFS\_v1 và BGL. Hai datasets này được chọn vì tính bổ sung (complementary): HDFS đại diện cho log data đơn giản và có cấu trúc, trong khi BGL đại diện cho log data phức tạp và đa dạng.

\subsection{HDFS\_v1: Hadoop Distributed File System}

\subsubsection{Nguồn gốc}

HDFS\_v1 được thu thập từ một Hadoop cluster 203 nodes chạy MapReduce workloads trong 38 giờ liên tục (năm 2008). Dataset này được tạo ra bởi Xu et al. \cite{xu2009detecting} trong nghiên cứu tiên phong về log mining.

\subsubsection{Đặc điểm}

\begin{table}[h]
\centering
\caption{Thống kê HDFS\_v1 dataset}
\label{tab:hdfs-stats}
\begin{tabular}{lr}
\toprule
\textbf{Characteristic} & \textbf{Value} \\
\midrule
Total log lines & 11,175,629 \\
Total blocks (sessions) & 575,061 \\
Unique templates & 29 \\
Average sequence length & 19.4 events \\
Normal blocks & 558,223 (97.07\%) \\
Anomaly blocks & 16,838 (2.93\%) \\
Imbalance ratio & 1:33 \\
\bottomrule
\end{tabular}
\end{table}

\subsubsection{Labeling strategy}

HDFS\_v1 sử dụng \textbf{fault injection} để tạo anomalies:

\begin{enumerate}
    \item Chạy healthy workload → Thu thập normal logs
    \item Inject faults:
    \begin{itemize}
        \item Kill DataNode processes
        \item Corrupt disk blocks
        \item Network partitions
        \item Disk full errors
    \end{itemize}
    \item Thu thập logs during failures → Label là anomalies
\end{enumerate}

\textbf{Advantages:}
\begin{itemize}
    \item Clean labels (biết chính xác khi nào anomaly xảy ra)
    \item Reproducible (có thể inject lại cùng faults)
    \item Controlled (biết failure modes)
\end{itemize}

\textbf{Disadvantages:}
\begin{itemize}
    \item Synthetic (có thể không reflect real production failures)
    \item Limited diversity (chỉ các failure modes được inject)
\end{itemize}

\subsubsection{Session grouping}

HDFS logs được nhóm theo \textbf{Block ID}. Mỗi block trong HDFS có một workflow rõ ràng:

\begin{verbatim}
Block blk_123 workflow:
  1. Client requests block write
  2. DataNode receives block
  3. Write to local disk
  4. Replicate to other DataNodes (typically 3 copies)
  5. Send acknowledgment to NameNode
\end{verbatim}

Tất cả logs liên quan đến cùng một Block ID được nhóm thành một \textbf{session}. Anomaly detection được thực hiện ở mức session (không phải individual log line).

\textbf{Label:} Session là anomalous nếu BẤT KỲ log nào trong session là anomalous.

\subsection{BGL: BlueGene/L Supercomputer}

\subsubsection{Nguồn gốc}

BGL được thu thập từ BlueGene/L supercomputer tại Lawrence Livermore National Lab trong 214 ngày (năm 2005). Supercomputer này có 128K processors và chạy các scientific computing workloads.

\subsubsection{Đặc điểm}

\begin{table}[h]
\centering
\caption{Thống kê BGL dataset}
\label{tab:bgl-stats}
\begin{tabular}{lr}
\toprule
\textbf{Characteristic} & \textbf{Value} \\
\midrule
Total log lines & 4,747,963 (raw) \\
Parsed successfully & 4,365,033 (91.9\%) \\
Unique templates & 1,775 \\
Normal logs & 3,837,629 (87.92\%) \\
Anomaly logs & 527,404 (12.08\%) \\
Imbalance ratio & 1:7 \\
Collection period & 214 days \\
\bottomrule
\end{tabular}
\end{table}

\subsubsection{Labeling strategy}

BGL sử dụng \textbf{natural errors} từ production system. Labels được trích xuất từ log level field:

\begin{itemize}
    \item \textbf{Anomaly:} FATAL, SEVERE, ERROR
    \item \textbf{Normal:} WARNING, INFO, DEBUG
\end{itemize}

\textbf{Advantages:}
\begin{itemize}
    \item Realistic (actual production failures)
    \item Diverse (214 days → nhiều failure modes)
    \item Unbiased (không có synthetic patterns)
\end{itemize}

\textbf{Disadvantages:}
\begin{itemize}
    \item Noisy labels (log level không hoàn toàn tương đương anomaly)
    \item Ground truth unclear (không có oracle)
    \item Higher anomaly rate (12\% vs expected 0.4\% trong literature - có thể do label extraction aggressive)
\end{itemize}

\subsubsection{Session grouping}

Khác với HDFS, BGL không có natural session ID. Do đó, sử dụng \textbf{time-window based grouping}:

\begin{itemize}
    \item \textbf{Window size:} 20 seconds (theo DeepLog \cite{du2017deeplog})
    \item \textbf{Sliding:} 10 seconds overlap (50\%)
    \item \textbf{Label:} Window là anomalous nếu BẤT KỲ log nào trong window là anomalous
\end{itemize}

\textbf{Rationale:} 20 seconds là khoảng thời gian đủ dài để capture context nhưng không quá dài để dilute anomaly signal.

\subsection{So sánh HDFS\_v1 vs BGL}

Hai datasets bổ sung cho nhau theo nhiều chiều:

\begin{table}[h]
\centering
\caption{So sánh HDFS\_v1 và BGL}
\label{tab:dataset-comparison}
\begin{tabular}{lll}
\toprule
\textbf{Aspect} & \textbf{HDFS\_v1} & \textbf{BGL} \\
\midrule
\textbf{Complexity} & Simple (29 templates) & Complex (1,775) \\
\textbf{Realism} & Synthetic (injection) & Natural (production) \\
\textbf{Structure} & Workflow (Block-based) & Event (time-based) \\
\textbf{Imbalance} & Severe (1:33) & Moderate (1:7) \\
\textbf{Duration} & Short (38 hours) & Long (214 days) \\
\textbf{Labeling} & Clean (injected) & Noisy (heuristic) \\
\bottomrule
\end{tabular}
\end{table}

\textbf{Complementarity:}
\begin{itemize}
    \item HDFS test methods trên simple, structured logs
    \item BGL test methods trên complex, diverse logs
    \item Together: Comprehensive coverage của LogAD challenges
\end{itemize}

\section{Phát hiện bất thường trên log}
\label{sec:log-anomaly-detection}

\subsection{Định nghĩa bài toán}

\textbf{Log-based Anomaly Detection (LogAD)} là bài toán phân loại (classification): Cho một log session (hoặc window), xác định nó là \textbf{Normal} hay \textbf{Anomaly}.

\textbf{Formal definition:}

\begin{itemize}
    \item \textbf{Input:} Log session $S = \{e_1, e_2, ..., e_n\}$ với $e_i$ là log events
    \item \textbf{Output:} Label $y \in \{0, 1\}$ với 0 = Normal, 1 = Anomaly
    \item \textbf{Goal:} Học function $f: S \rightarrow y$ từ training data
\end{itemize}

\subsection{Các hướng tiếp cận}

Dựa trên mức độ supervision, LogAD methods được phân loại thành:

\subsubsection{Supervised Learning}

Yêu cầu labeled data (cả normal và anomaly) cho training.

\textbf{Examples:}
\begin{itemize}
    \item LightAD Decision Tree, SLFN (khi train với labels)
    \item LogBERT fine-tuned với labeled data
\end{itemize}

\textbf{Advantages:}
\begin{itemize}
    \item High accuracy (nếu labels chính xác)
    \item Có thể học specific anomaly patterns
\end{itemize}

\textbf{Disadvantages:}
\begin{itemize}
    \item Yêu cầu nhiều labeled data (expensive)
    \item Không detect được novel anomaly types chưa thấy trong training
\end{itemize}

\subsubsection{Semi-supervised Learning}

Train trên mostly normal data với một ít labeled anomalies.

\textbf{Examples:}
\begin{itemize}
    \item LogRAG: DeepSVDD train trên normal data, sau đó dùng few-shot LLM
    \item One-Class SVM
\end{itemize}

\textbf{Advantages:}
\begin{itemize}
    \item Ít labels hơn supervised
    \item Có thể detect một số novel anomalies
\end{itemize}

\textbf{Disadvantages:}
\begin{itemize}
    \item Vẫn cần một số labeled data
    \item Performance phụ thuộc vào quality của normal data
\end{itemize}

\subsubsection{Unsupervised Learning}

Không cần labels, chỉ dựa vào distribution của data.

\textbf{Examples:}
\begin{itemize}
    \item LightAD KNN (distance-based)
    \item PCA-based methods
    \item Clustering-based methods
\end{itemize}

\textbf{Advantages:}
\begin{itemize}
    \item Không cần labels (practical)
    \item Có thể detect novel anomalies
\end{itemize}

\textbf{Disadvantages:}
\begin{itemize}
    \item Lower accuracy
    \item Khó tune threshold
    \item High false positive rate
\end{itemize}

\subsection{Evaluation metrics}

Do tính imbalanced của LogAD, metrics phải được chọn cẩn thận:

\subsubsection{Precision, Recall, F1-score}

\begin{align}
Precision &= \frac{TP}{TP + FP} \\
Recall &= \frac{TP}{TP + FN} \\
F1 &= \frac{2 \times Precision \times Recall}{Precision + Recall}
\end{align}

Trong đó:
\begin{itemize}
    \item TP (True Positive): Anomaly predicted đúng
    \item FP (False Positive): Normal predicted nhầm là anomaly
    \item FN (False Negative): Anomaly predicted nhầm là normal
    \item TN (True Negative): Normal predicted đúng
\end{itemize}

\textbf{F1-score} là harmonic mean của Precision và Recall, cân bằng giữa hai metrics này. F1 được sử dụng làm \textbf{primary metric} trong nghiên cứu này.

\subsubsection{PR-AUC (Precision-Recall Area Under Curve)}

PR-AUC là diện tích dưới đường cong Precision-Recall. Metric này:
\begin{itemize}
    \item Threshold-independent (không phụ thuộc vào threshold cụ thể)
    \item Phù hợp cho imbalanced data (hơn ROC-AUC)
    \item Tổng hợp performance across all thresholds
\end{itemize}

\textbf{Range:} [0, 1], càng cao càng tốt.

\subsubsection{False Positive Rate (FPR)}

\begin{equation}
FPR = \frac{FP}{FP + TN}
\end{equation}

FPR đo tỷ lệ false alarms. Trong production, high FPR gây ra:
\begin{itemize}
    \item Alarm fatigue (SREs ignore alerts)
    \item Wasted investigation time
    \item Reduced trust in system
\end{itemize}

\textbf{Goal:} FPR càng thấp càng tốt (thường <5\%).

\subsubsection{Why not Accuracy?}

Accuracy $= \frac{TP + TN}{TP + TN + FP + FN}$ là misleading cho imbalanced data.

\textbf{Example (BGL with 12\% anomaly):}
\begin{itemize}
    \item Predict ALL logs as "Normal" → Accuracy = 88\%!
    \item But Recall = 0\% (miss all anomalies)
\end{itemize}

Do đó, accuracy KHÔNG được sử dụng trong nghiên cứu này.

\section{Tổng quan các phương pháp}
\label{sec:methods-overview}

Phần này giới thiệu tổng quan ba phương pháp được đánh giá trong luận văn. Chi tiết đầy đủ về implementation và reproduction được trình bày trong Chương 3.

\subsection{LightAD: Classical Machine Learning}

LightAD \cite{yu2024lightad} đại diện cho classical ML approach với ba models:

\subsubsection{K-Nearest Neighbors (KNN)}

\textbf{Idea:} Anomaly là samples "xa" với majority samples.

\textbf{Algorithm:}
\begin{enumerate}
    \item Represent mỗi session bằng event count vector (bag-of-events)
    \item For test sample, tìm $k$ nearest neighbors trong training set
    \item Nếu majority neighbors là anomaly → Predict anomaly
    \item Hoặc dùng average distance > threshold
\end{enumerate}

\textbf{Complexity:} O(n) inference (cần scan all training samples)

\textbf{Explainability:} Có thể show nearest neighbors

\subsubsection{Decision Tree (DT)}

\textbf{Idea:} Học if-then rules để phân loại.

\textbf{Advantages:}
\begin{itemize}
    \item Fast inference: O(log n)
    \item Highly interpretable (rules dễ hiểu)
    \item No assumption về distribution
\end{itemize}

\textbf{Explainability:} Decision path cung cấp rules rõ ràng

\subsubsection{Single Layer Feedforward Network (SLFN)}

\textbf{Idea:} Shallow neural network với 1 hidden layer.

\textbf{Architecture:}
\begin{verbatim}
Input (event counts) → Hidden (100 neurons, ReLU) → Output (2 classes, Softmax)
\end{verbatim}

\textbf{Advantages:}
\begin{itemize}
    \item More expressive than linear models
    \item Faster than deep networks
\end{itemize}

\textbf{Explainability:} Feature importance via gradients

\subsection{LogFormer: Pre-trained Transformer}

LogFormer \cite{guo2024logformer} áp dụng pre-training + adapter tuning paradigm.

\subsubsection{Architecture}

\textbf{Phase 1: Pre-training (Unsupervised)}
\begin{itemize}
    \item Train Transformer Encoder trên large unlabeled log corpus (50M logs)
    \item Task: Masked Language Modeling (MLM) - predict masked log events
    \item Goal: Learn general log patterns
\end{itemize}

\textbf{Phase 2: Adapter Tuning (Supervised)}
\begin{itemize}
    \item Freeze pre-trained weights
    \item Insert lightweight adapter modules (~3M params)
    \item Fine-tune chỉ adapters trên target dataset
    \item Task: Binary classification (normal/anomaly)
\end{itemize}

\subsubsection{Advantages}

\begin{itemize}
    \item Transfer learning: Leverage knowledge từ nhiều datasets
    \item Parameter-efficient: Chỉ tune 3M params thay vì 110M
    \item Fast fine-tuning: 10-30 mins thay vì 4-8 hours
    \item High accuracy: State-of-the-art trên nhiều benchmarks
\end{itemize}

\subsubsection{Explainability}

Attention weights cho thấy model focus vào log events nào, nhưng khó interpret hơn Decision Tree rules.

\subsection{LogRAG: LLM with Retrieval-Augmented Generation}

LogRAG \cite{jiang2024lograg} kết hợp DeepSVDD + RAG + LLM.

\subsubsection{Architecture}

\textbf{Phase 1: Anomaly Candidate Detection (DeepSVDD)}
\begin{itemize}
    \item Train DeepSVDD (deep one-class SVM) trên mostly normal data
    \item Learn hypersphere bao quanh normal samples
    \item Output: Anomaly scores cho all samples
    \item Select top-k samples với high scores làm candidates
\end{itemize}

\textbf{Phase 2: RAG-enhanced LLM Reasoning}
\begin{itemize}
    \item For mỗi candidate:
    \begin{enumerate}
        \item \textbf{Retrieval:} Tìm k=5 similar normal examples + k=5 similar anomaly examples (dùng FAISS)
        \item \textbf{Augmentation:} Construct prompt với retrieved examples + query sample
        \item \textbf{Generation:} LLM (GPT-3.5/GPT-4) generates reasoning + prediction
    \end{enumerate}
\end{itemize}

\subsubsection{Advantages}

\begin{itemize}
    \item High accuracy: State-of-the-art F1 scores
    \item Explainability: Natural language reasoning
    \item Few-shot: Chỉ cần ít labeled examples (semi-supervised)
    \item Semantic understanding: LLM hiểu ý nghĩa log messages
\end{itemize}

\subsubsection{Disadvantages}

\begin{itemize}
    \item High latency: 50-100ms per sample (LLM slow)
    \item API cost: \$2-5 per 1K samples (GPT-3.5)
    \item Hallucination risk: LLM có thể generate plausible but incorrect explanations
\end{itemize}

\section{Tóm tắt chương}

Chương này đã trình bày các kiến thức nền tảng:

\begin{itemize}
    \item \textbf{Dữ liệu log:} 5 đặc điểm chính (semi-structured, non-stationary, high-velocity, imbalanced, contextual)
    \item \textbf{Log parsing:} Drain algorithm để extract templates
    \item \textbf{Datasets:} HDFS\_v1 (simple, controlled) và BGL (complex, natural) - complementary
    \item \textbf{LogAD problem:} Định nghĩa, supervised/semi-supervised/unsupervised approaches
    \item \textbf{Metrics:} Precision, Recall, F1, PR-AUC, FPR (không dùng accuracy)
    \item \textbf{Phương pháp:} LightAD (classical ML), LogFormer (Transformer), LogRAG (LLM+RAG)
\end{itemize}

Chương tiếp theo (Chương 3) sẽ trình bày chi tiết phương pháp nghiên cứu và giao thức thực nghiệm.

% \chapter{Phương pháp nghiên cứu}
\label{chap:methodology}

% ============================================
% CHAPTER 3: METHODOLOGY
% ============================================
% Task 2.11-2.14: Viết Chapter 3 (15-20 pages)
% Trình bày chi tiết experimental protocol, RQs, pipeline, evaluation

\section{Tổng quan phương pháp}
\label{sec:methodology-overview}

Luận văn này thực hiện một \textit{empirical benchmarking study} (nghiên cứu đánh giá thực nghiệm) nhằm so sánh ba phương pháp phát hiện bất thường trên log đại diện cho ba thế hệ machine learning: Classical ML (LightAD), Deep Learning (LogFormer), và LLM (LogRAG). 

\textbf{Loại nghiên cứu:} Quantitative empirical study tuân theo ACM SIGSOFT Empirical Standards \cite{acm_empirical_standards}.

\subsection{Mục tiêu nghiên cứu}
\label{subsec:research-goals}

\textbf{Mục tiêu chính:} Trả lời câu hỏi: "\textit{Phương pháp nào là lựa chọn tốt nhất cho log anomaly detection trong thực tế, xét trên nhiều tiêu chí: hiệu quả, chi phí, và khả năng giải thích?}"

\textbf{Khác biệt so với các nghiên cứu trước:}
\begin{itemize}
    \item \textbf{Cross-generational comparison:} Đầu tiên so sánh hệ thống ba thế hệ ML (Classical, DL, LLM) trên cùng datasets
    \item \textbf{Multi-dimensional evaluation:} Không chỉ accuracy, mà cả computational cost, explainability, error patterns
    \item \textbf{Fair protocol:} Cùng data splits, cùng hardware, cùng metrics
    \item \textbf{Practical guidelines:} Khuyến nghị triển khai dựa trên kết quả thực nghiệm
\end{itemize}

\subsection{Tiếp cận theo RQ (Research Questions)}
\label{subsec:rq-driven}

Nghiên cứu được tổ chức xung quanh 4 Research Questions (RQ), từ hiệu quả đến chi phí, trade-offs, và explainability. Mỗi RQ có metrics và phương pháp đánh giá cụ thể, được trình bày chi tiết trong Section \ref{sec:research-questions}.

\section{Câu hỏi nghiên cứu}
\label{sec:research-questions}

\subsection{RQ1: Hiệu quả phát hiện bất thường}
\label{subsec:rq1}

\textbf{Research Question 1:}

\begin{quote}
\textit{Hiệu quả phát hiện bất thường (detection effectiveness) của LightAD, LogFormer, và LogRAG là bao nhiêu khi đánh giá trên HDFS\_v1 và BGL datasets?}
\end{quote}

\subsubsection{Động lực}

Hiệu quả phát hiện là tiêu chí quan trọng nhất. Một phương pháp không phát hiện được anomaly (low Recall) hoặc tạo ra quá nhiều false alarms (low Precision) đều không khả thi trong production.

\subsubsection{Metrics}

Chúng tôi đo lường hiệu quả qua 5 metrics:

\begin{enumerate}
    \item \textbf{Precision (P):} Tỷ lệ dự đoán anomaly chính xác
    \[
    P = \frac{TP}{TP + FP}
    \]
    \textit{Ý nghĩa:} P cao = ít false alarms = ít "alert fatigue"
    
    \item \textbf{Recall (R):} Tỷ lệ anomaly được phát hiện
    \[
    R = \frac{TP}{TP + FN}
    \]
    \textit{Ý nghĩa:} R cao = ít anomaly bị bỏ sót = ít incidents
    
    \item \textbf{F1-Score (F1):} Harmonic mean của P và R
    \[
    F1 = \frac{2 \cdot P \cdot R}{P + R}
    \]
    \textit{Ý nghĩa:} F1 cân bằng Precision và Recall, metric chính để ranking methods
    
    \item \textbf{PR-AUC:} Area under Precision-Recall curve
    
    \textit{Ý nghĩa:} Đánh giá performance across all thresholds, robust với imbalanced data
    
    \item \textbf{FPR:} False Positive Rate
    \[
    FPR = \frac{FP}{FP + TN}
    \]
    \textit{Ý nghĩa:} FPR < 1\% thường được chấp nhận trong industry (target)
\end{enumerate}

\subsubsection{Kỳ vọng}

\textbf{Hypothesis:} Deep Learning và LLM methods (LogFormer, LogRAG) sẽ có F1 và PR-AUC cao hơn Classical ML (LightAD), nhưng sự khác biệt có thể không lớn trên datasets đơn giản như HDFS\_v1.

\textbf{Insight mong đợi:} Xác định phương pháp nào có accuracy cao nhất, và liệu classical ML có "competitive" hay không (theo tuyên bố của LightAD paper \cite{yu2024lightad}).

\subsection{RQ2: Chi phí tính toán}
\label{subsec:rq2}

\textbf{Research Question 2:}

\begin{quote}
\textit{Chi phí tính toán (computational cost) của mỗi phương pháp là bao nhiêu, bao gồm training time, inference latency, memory usage, và throughput?}
\end{quote}

\subsubsection{Động lực}

Accuracy không phải là tiêu chí duy nhất. Một phương pháp với F1 = 0.98 nhưng inference time 10 seconds/sample không thể triển khai trong real-time system. Chi phí cũng bao gồm API cost cho LLM-based methods.

\subsubsection{Metrics}

\begin{enumerate}
    \item \textbf{Training Time:} Thời gian huấn luyện model từ đầu (wall-clock time)
    
    \textit{Measured:} Seconds (s) hoặc Hours (h)
    
    \textit{Ý nghĩa:} Quan trọng khi cần retrain thường xuyên (concept drift)
    
    \item \textbf{Inference Latency:} Thời gian dự đoán một sample
    
    \textit{Measured:} Milliseconds (ms), report mean, p95, p99
    
    \textit{Ý nghĩa:} Critical cho real-time systems (target: < 100ms/sample)
    
    \item \textbf{Memory Usage:}
    \begin{itemize}
        \item \textbf{Peak RAM:} System memory (GB)
        \item \textbf{Peak VRAM:} GPU memory (GB)
    \end{itemize}
    
    \textit{Ý nghĩa:} Quan trọng cho resource-constrained environments
    
    \item \textbf{Throughput:} Số samples xử lý mỗi giây
    
    \textit{Measured:} Samples/second
    
    \textit{Ý nghĩa:} Quan trọng cho high-velocity log streams (millions logs/hour)
    
    \item \textbf{API Cost (LogRAG only):} Chi phí OpenAI API
    
    \textit{Measured:} USD per 1K samples
    
    \textit{Ý nghĩa:} Operational cost cho LLM-based methods
\end{enumerate}

\subsubsection{Phương pháp đo}

\begin{itemize}
    \item \textbf{Hardware cố định:} NVIDIA RTX 3090 (24GB VRAM), Intel Core i9, 32GB RAM
    \item \textbf{Monitoring tools:} 
        \begin{itemize}
            \item \texttt{nvidia-smi} cho GPU metrics
            \item \texttt{psutil} (Python) cho RAM
            \item \texttt{time.perf\_counter()} cho wall-clock time
        \end{itemize}
    \item \textbf{Repeated runs:} Mỗi measurement chạy 5 lần, report mean $\pm$ std
\end{itemize}

\subsubsection{Kỳ vọng}

\textbf{Hypothesis:} 
\begin{itemize}
    \item Training time: LightAD < LogFormer < LogRAG (không có pre-training overhead)
    \item Inference latency: LightAD (DT) < LogFormer < LogRAG (API calls)
    \item Memory: LightAD < LogFormer < LogRAG
    \item API cost: LogRAG đắt nhất (GPT-3.5: \$0.002/1K tokens)
\end{itemize}

\textbf{Insight mong đợi:} Quantify chi phí của deep learning và LLM so với classical ML. Xác định liệu classical ML có "efficiency advantage" đáng kể hay không.

\subsection{RQ3: Đánh đổi hiệu năng-chi phí}
\label{subsec:rq3}

\textbf{Research Question 3:}

\begin{quote}
\textit{Đánh đổi (trade-offs) giữa detection effectiveness (RQ1) và computational cost (RQ2) là gì? Phương pháp nào nằm trên Pareto frontier?}
\end{quote}

\subsubsection{Động lực}

Trong thực tế, lựa chọn phương pháp phụ thuộc vào deployment constraints:
\begin{itemize}
    \item High-stakes production (banking, healthcare): Ưu tiên accuracy
    \item Real-time systems (network monitoring): Ưu tiên latency
    \item Resource-constrained (edge devices): Ưu tiên memory
\end{itemize}

RQ3 cung cấp trade-off analysis để hỗ trợ decision-making.

\subsubsection{Phương pháp phân tích}

\textbf{Pareto Frontier Analysis:}

\begin{enumerate}
    \item Vẽ scatter plots: 
    \begin{itemize}
        \item X-axis: Inference latency (ms)
        \item Y-axis: F1-Score
    \end{itemize}
    
    \item Xác định Pareto-optimal points: Phương pháp $A$ dominates $B$ nếu:
    \[
    F1_A \geq F1_B \quad \text{AND} \quad \text{Latency}_A \leq \text{Latency}_B
    \]
    với ít nhất một inequality strict.
    
    \item Phương pháp trên Pareto frontier là optimal trade-off (không thể improve một metric mà không sacrifice metric kia)
\end{enumerate}

\textbf{Multi-dimensional comparison:}
\begin{itemize}
    \item Radar charts: So sánh 3 methods trên 5 dimensions (F1, Speed, Memory, Cost, Explainability)
    \item Normalize mỗi dimension về 0-1 scale
\end{itemize}

\subsubsection{Kỳ vọng}

\textbf{Hypothesis:}
\begin{itemize}
    \item LightAD-DT có thể là Pareto-optimal: Good F1 + Very low latency
    \item LogFormer: Higher F1 nhưng moderate latency
    \item LogRAG: Highest F1 nhưng highest latency và cost
\end{itemize}

\textbf{Insight mong đợi:} Decision tree để chọn phương pháp dựa trên use case constraints.

\subsection{RQ4: Khả năng giải thích và phân tích lỗi}
\label{subsec:rq4}

\textbf{Research Question 4:}

\begin{quote}
\textit{Các error patterns (FP/FN) của mỗi phương pháp là gì? Khả năng giải thích (explainability) của predictions như thế nào?}
\end{quote}

\subsubsection{Động lực}

\textbf{Phần (a) - Error Analysis:} Hiểu \textit{tại sao} một phương pháp fail giúp:
\begin{itemize}
    \item Identify weaknesses (e.g., poor on rare templates)
    \item Guide method selection (e.g., avoid method X nếu dataset có characteristic Y)
    \item Future improvements (e.g., hybrid approaches)
\end{itemize}

\textbf{Phần (b) - Explainability:} Trong production, SREs cần biết \textit{tại sao} một log được classify as anomaly. Black-box predictions không được trust.

\subsubsection{Phương pháp}

\textbf{(a) Error Analysis:}

\begin{enumerate}
    \item \textbf{Extract errors:} Từ test set, lấy 50-100 False Positives (FP) và 50-100 False Negatives (FN) của mỗi phương pháp.
    
    \item \textbf{Build error taxonomy:} Classify errors vào các categories:
    \begin{itemize}
        \item \textit{Unseen template:} Log template chưa thấy trong training (rare event)
        \item \textit{Short sequence:} Sequence quá ngắn, thiếu context
        \item \textit{Noise:} Parsing errors, malformed logs
        \item \textit{Class imbalance:} Minority class bị underrepresented
        \item \textit{Semantic ambiguity:} Log trông bình thường nhưng thực tế bất thường (hoặc ngược lại)
    \end{itemize}
    
    \item \textbf{Quantify distribution:} Tính phần trăm mỗi error category cho mỗi phương pháp.
    
    \item \textbf{Case studies:} Phân tích chi tiết 2-3 FP và 2-3 FN examples mỗi phương pháp.
\end{enumerate}

\textbf{(b) Explainability Analysis:}

\begin{enumerate}
    \item \textbf{Extract explanations:}
    \begin{itemize}
        \item \textbf{LightAD-KNN:} Show k nearest neighbors
        \item \textbf{LightAD-DT:} Show decision path (which rules applied)
        \item \textbf{LightAD-SLFN:} Feature importance (top features)
        \item \textbf{LogFormer:} Attention weights (which log events attended)
        \item \textbf{LogRAG:} Retrieved examples + LLM reasoning chain
    \end{itemize}
    
    \item \textbf{Assess explainability criteria:}
    \begin{itemize}
        \item \textit{Actionability:} Liệu SRE có thể sử dụng explanation để debug? (1-5 scale)
        \item \textit{Clarity:} Explanation có dễ hiểu với non-ML expert? (1-5 scale)
        \item \textit{Correctness:} Explanation có align với root cause thực tế? (1-5 scale)
    \end{itemize}
    
    \item \textbf{Visualizations:}
    \begin{itemize}
        \item Decision tree diagram (LightAD-DT)
        \item Attention heatmap (LogFormer)
        \item Retrieved examples + reasoning chain (LogRAG)
    \end{itemize}
\end{enumerate}

\subsubsection{Kỳ vọng}

\textbf{Hypothesis:}
\begin{itemize}
    \item \textbf{Error patterns:} 
        \begin{itemize}
            \item LightAD: Nhiều errors trên rare templates (limited training data)
            \item LogFormer: Ít errors hơn do pre-training (learned general patterns)
            \item LogRAG: Ít FN (high Recall) nhưng có thể nhiều FP (over-sensitive)
        \end{itemize}
    \item \textbf{Explainability:}
        \begin{itemize}
            \item LightAD-DT: Best (clear if-then rules)
            \item LogRAG: Best for non-ML users (natural language)
            \item LogFormer: Moderate (attention weights need expertise)
        \end{itemize}
\end{itemize}

\textbf{Insight mong đợi:} 
\begin{itemize}
    \item Identify which method is best for explainability-critical applications
    \item Common failure modes across methods (guide future work)
\end{itemize}

\section{Pipeline thực nghiệm}
\label{sec:experimental-pipeline}

\subsection{Quy trình end-to-end}
\label{subsec:pipeline-overview}

Hình \ref{fig:pipeline} minh họa pipeline thực nghiệm 5 bước:

\begin{figure}[ht]
\centering
\begin{tikzpicture}[
    node distance=1.5cm,
    box/.style={rectangle, draw, fill=blue!20, text width=3cm, text centered, minimum height=1cm},
    arrow/.style={->, >=stealth, thick}
]
    \node[box] (collect) {1. Data Collection};
    \node[box, below of=collect] (parse) {2. Log Parsing};
    \node[box, below of=parse] (split) {3. Data Splitting};
    \node[box, below of=split] (train) {4. Model Training};
    \node[box, below of=train] (eval) {5. Evaluation};
    
    \draw[arrow] (collect) -- (parse);
    \draw[arrow] (parse) -- (split);
    \draw[arrow] (split) -- (train);
    \draw[arrow] (train) -- (eval);
\end{tikzpicture}
\caption{Pipeline thực nghiệm 5 bước}
\label{fig:pipeline}
\end{figure}

\subsubsection{Step 1: Data Collection}

\textbf{Datasets:} HDFS\_v1 và BGL từ LogHub \cite{he2020loghub}.

\textbf{Characteristics:} Đã trình bày chi tiết trong Section \ref{subsec:datasets} (Chapter 2).

\textbf{Download:}
\begin{itemize}
    \item \textbf{HDFS\_v1:} \url{https://zenodo.org/record/8196385}
    \item \textbf{BGL:} \url{https://zenodo.org/record/8196385}
\end{itemize}

\textbf{Verification:} Checksum verification để đảm bảo data integrity.

\subsubsection{Step 2: Log Parsing}

\textbf{Tool:} Drain3 \cite{he2017drain} - state-of-art log parser.

\textbf{Purpose:} Chuyển đổi raw logs thành structured format:
\begin{itemize}
    \item \textbf{Input:} Raw log line (semi-structured text)
    \item \textbf{Output:} 
        \begin{itemize}
            \item \textit{Template ID:} Integer representing log event type
            \item \textit{Template:} Template string (với <*> cho parameters)
            \item \textit{Parameters:} Extracted variable values
        \end{itemize}
\end{itemize}

\textbf{Example:}
\begin{verbatim}
Input:  "2024-01-15 10:23:45 ERROR Connection to 192.168.1.1 failed"
Output: {
    "template_id": 42,
    "template": "ERROR Connection to <*> failed",
    "parameters": ["192.168.1.1"],
    "timestamp": "2024-01-15 10:23:45"
}
\end{verbatim}

\textbf{Configuration:}
\begin{itemize}
    \item \textbf{Depth:} 4 (Drain tree depth)
    \item \textbf{Similarity threshold:} 0.5
    \item \textbf{Max children:} 100
\end{itemize}

\textbf{Output:}
\begin{itemize}
    \item \texttt{hdfs\_structured.jsonl} (575K blocks)
    \item \texttt{bgl\_structured.jsonl} (4.3M log lines)
    \item Template dictionaries với frequencies
\end{itemize}

\subsubsection{Step 3: Data Splitting}

\textbf{Ratio:} 80\% training, 10\% validation, 10\% test

\textbf{Strategy:} Khác nhau cho mỗi dataset để tránh data leakage và ensure temporal validity.

\paragraph{HDFS\_v1: Block-based Splitting}

\begin{itemize}
    \item \textbf{Grouping:} Mỗi block ID là một session
    \item \textbf{Label:} Block anomalous nếu bất kỳ log nào trong block là anomalous
    \item \textbf{Splitting:} Stratified split theo label (giữ anomaly ratio)
    \item \textbf{Rationale:} Tránh leakage - logs từ cùng block không được tách ra train/test
\end{itemize}

\textbf{Algorithm:}
\begin{verbatim}
1. Group logs by BlockId
2. Aggregate label: Block anomalous if max(log labels) = 1
3. Stratified split by block label:
   - Train: 80% blocks (stratified)
   - Val: 10% blocks (stratified)
   - Test: 10% blocks (stratified)
\end{verbatim}

\paragraph{BGL: Chronological Splitting}

\begin{itemize}
    \item \textbf{Grouping:} Sliding window (20 consecutive logs)
    \item \textbf{Label:} Window anomalous nếu bất kỳ log nào bất thường
    \item \textbf{Splitting:} Chronological (first 80\% → train, next 10\% → val, last 10\% → test)
    \item \textbf{Rationale:} Preserve temporal order, test on future data (realistic deployment)
\end{itemize}

\textbf{Algorithm:}
\begin{verbatim}
1. Sort logs by timestamp
2. Create sliding windows (size=20, step=10)
3. Label window: Anomalous if any(log labels) = 1
4. Chronological split:
   - Train: Windows 0 to 80%
   - Val: Windows 80% to 90%
   - Test: Windows 90% to 100%
\end{verbatim}

\textbf{Verification:}
\begin{itemize}
    \item Check no overlap between train/val/test
    \item Verify anomaly rates similar across splits (±2\%)
    \item Document split statistics
\end{itemize}

\subsubsection{Step 4: Model Training}

\textbf{Seed:} Fixed random seed = 42 cho reproducibility

\textbf{Hyperparameters:}
\begin{itemize}
    \item \textbf{Start:} Paper defaults cho mỗi phương pháp
    \item \textbf{Tuning:} Chỉ tune trên validation set
    \item \textbf{Rule:} \textbf{NEVER} touch test set (no peeking!)
\end{itemize}

\paragraph{LightAD Training}

\begin{itemize}
    \item \textbf{Feature representation:} Event count vector
    \item \textbf{Models:}
        \begin{itemize}
            \item KNN: k=5 (default from paper)
            \item Decision Tree: max\_depth=10, min\_samples\_split=20
            \item SLFN: hidden\_units=50, epochs=100, lr=0.001
        \end{itemize}
    \item \textbf{Training time:} Minutes (< 10 min expected)
\end{itemize}

\paragraph{LogFormer Training}

\begin{itemize}
    \item \textbf{Phase 1: Pre-training (optional - use pretrained BigLog model)}
        \begin{itemize}
            \item Model: BERT-base architecture
            \item Corpus: BigLog (100M logs)
            \item Task: Masked Language Modeling
        \end{itemize}
    \item \textbf{Phase 2: Adapter Tuning}
        \begin{itemize}
            \item Adapter config: bottleneck dim = 64
            \item Epochs: 10 (tune on val set)
            \item Batch size: 32
            \item Learning rate: 1e-4
        \end{itemize}
    \item \textbf{Training time:} Hours (4-8 hours on RTX 3090)
\end{itemize}

\paragraph{LogRAG Training}

\begin{itemize}
    \item \textbf{Phase 1: DeepSVDD}
        \begin{itemize}
            \item Encoder: 3-layer MLP [128, 64, 32]
            \item Hypersphere center: Initialized from mean of normal samples
            \item Epochs: 50
            \item Learning rate: 0.0001
        \end{itemize}
    \item \textbf{Phase 2: RAG Setup}
        \begin{itemize}
            \item Retriever: FAISS index (Flat L2)
            \item Top-k: 5 retrieved examples
            \item LLM: GPT-3.5-turbo (OpenAI API) hoặc Mistral-7B (local)
        \end{itemize}
    \item \textbf{Training time:} Hours (DeepSVDD: 2-4 hours, no LLM training)
\end{itemize}

\subsubsection{Step 5: Evaluation}

\textbf{Test set:} Chỉ sử dụng sau khi model đã fixed (no tuning on test!)

\textbf{Metrics:} Compute all metrics từ RQ1 (Precision, Recall, F1, PR-AUC, FPR)

\textbf{Statistical testing:} Wilcoxon signed-rank test cho pairwise comparisons (significant if p < 0.05)

\textbf{Repeated runs:} 5 runs với different seeds (42, 123, 456, 789, 2024), report mean ± std

\subsection{Kiểm soát công bằng}
\label{subsec:fairness-controls}

Để đảm bảo so sánh công bằng giữa 3 phương pháp, chúng tôi kiểm soát các factors sau:

\subsubsection{Phần cứng thống nhất}

\textbf{Hardware configuration:}
\begin{itemize}
    \item \textbf{GPU:} NVIDIA RTX 3090 (24GB VRAM)
    \item \textbf{CPU:} Intel Core i9-12900K
    \item \textbf{RAM:} 32GB DDR4
    \item \textbf{Storage:} 1TB NVMe SSD
    \item \textbf{OS:} Ubuntu 22.04 LTS
\end{itemize}

\textbf{Rationale:} Eliminate hardware confounding factor. Tất cả methods chạy trên cùng machine.

\subsubsection{Seed cố định}

\textbf{Random seed = 42} cho:
\begin{itemize}
    \item Data splitting (train/val/test)
    \item Model initialization
    \item Stochastic training procedures
\end{itemize}

\textbf{Rationale:} Reproducibility và eliminate randomness.

\textbf{Multiple runs:} Chạy 5 lần với 5 seeds khác nhau (42, 123, 456, 789, 2024) để estimate variance.

\subsubsection{Hyperparameters}

\textbf{Protocol:}
\begin{enumerate}
    \item \textbf{Start:} Paper defaults cho mỗi method
    \item \textbf{Tuning:} Nếu cần tune, chỉ tune trên validation set
    \item \textbf{Rule:} No peeking at test set
    \item \textbf{Document:} Ghi lại tất cả hyperparameters trong appendix
\end{enumerate}

\textbf{Rationale:} Fair comparison - no cherry-picking hyperparameters to favor one method.

\subsubsection{Evaluation protocol}

\textbf{Same for all methods:}
\begin{itemize}
    \item Same data splits (train/val/test)
    \item Same metrics (P/R/F1/PR-AUC/FPR)
    \item Same evaluation script (unified \texttt{evaluate.py})
    \item Same statistical tests
\end{itemize}

\textbf{Rationale:} Apples-to-apples comparison.

\section{Giao thức tái lập}
\label{sec:reproducibility}

Chúng tôi cam kết reproducibility theo ACM Artifact Review and Badging guidelines \cite{acm_artifact_badging}.

\subsection{Code release}

\textbf{GitHub repository:} Public repository chứa:
\begin{itemize}
    \item All scripts (parsing, training, evaluation, visualization)
    \item Configuration files (hyperparameters)
    \item Data processing code
    \item Jupyter notebooks cho analysis
\end{itemize}

\textbf{License:} MIT License (permissive open-source)

\subsection{Data availability}

\textbf{Datasets:} Links đến LogHub (HDFS\_v1, BGL)

\textbf{Processed data:} Parsed logs và splits sẽ được host trên Zenodo (DOI)

\textbf{Privacy:} Cả hai datasets đều public, không có privacy concerns.

\subsection{Environment}

\textbf{Docker container:} Dockerfile để recreate exact environment:
\begin{itemize}
    \item Python 3.8
    \item PyTorch 1.13.1
    \item Transformers 4.30.0
    \item All dependencies với pinned versions
\end{itemize}

\textbf{Quick start:}
\begin{verbatim}
docker pull <username>/logad-bench:v1.0
docker run --gpus all -v $(pwd):/workspace <image>
bash scripts/run_all_experiments.sh
\end{verbatim}

\subsection{Results package}

\textbf{Contents:}
\begin{itemize}
    \item Raw results (JSON files với tất cả metrics)
    \item Trained model checkpoints (nếu khả thi, exclude LogRAG do API)
    \item Analysis scripts (reproduce figures và tables)
    \item LaTeX source cho thesis chapters
\end{itemize}

\textbf{Goal:} Any researcher có thể reproduce mọi figure và table trong thesis.

\section{Threats to Validity}
\label{sec:threats-to-validity}

Theo ACM Empirical Standards, chúng tôi identify và mitigate các threats to validity:

\subsection{Internal Validity}

\textbf{Threat 1: Implementation bugs}

\textit{Mitigation:}
\begin{itemize}
    \item Sử dụng official implementations (LightAD, LogFormer, LogRAG GitHub repos)
    \item Unit tests cho data processing code
    \item Sanity checks (e.g., metrics in reasonable range)
\end{itemize}

\textbf{Threat 2: Hyperparameter tuning bias}

\textit{Mitigation:}
\begin{itemize}
    \item Start from paper defaults
    \item Document all tuning decisions
    \item Tune only on validation set, never test set
\end{itemize}

\textbf{Threat 3: Random variation}

\textit{Mitigation:}
\begin{itemize}
    \item Fixed random seeds
    \item Multiple runs (5 seeds)
    \item Report mean ± std
    \item Statistical significance testing
\end{itemize}

\subsection{External Validity}

\textbf{Threat 1: Limited datasets}

\textit{Discussion:} Chỉ 2 datasets (HDFS\_v1, BGL). Results có thể không generalize sang hệ thống khác.

\textit{Mitigation:}
\begin{itemize}
    \item Chọn 2 datasets representative: HDFS (simple, structured) vs BGL (complex, diverse)
    \item Acknowledge limitation trong discussion
    \item Future work: Extend to more datasets
\end{itemize}

\textbf{Threat 2: Specific versions}

\textit{Discussion:} Kết quả dựa trên specific versions của LightAD, LogFormer, LogRAG. Newer versions có thể khác.

\textit{Mitigation:}
\begin{itemize}
    \item Document exact versions và commit hashes
    \item Pin dependencies trong Docker
    \item Acknowledge limitation
\end{itemize}

\subsection{Construct Validity}

\textbf{Threat 1: Metric choices}

\textit{Discussion:} Liệu F1, PR-AUC có phải best metrics cho LogAD?

\textit{Mitigation:}
\begin{itemize}
    \item F1 và PR-AUC widely accepted trong LogAD literature
    \item Report multiple metrics (5 total) để triangulate
    \item RQ4 provides qualitative analysis (error patterns)
\end{itemize}

\textbf{Threat 2: Explainability assessment}

\textit{Discussion:} Actionability, Clarity, Correctness là subjective (assessed by authors).

\textit{Mitigation:}
\begin{itemize}
    \item Provide concrete examples (case studies)
    \item Document assessment criteria clearly
    \item Acknowledge subjectivity
    \item Future work: User study với SREs
\end{itemize}

\section{Kết luận}
\label{sec:methodology-conclusion}

Chương này đã trình bày:
\begin{enumerate}
    \item \textbf{4 Research Questions:} RQ1 (Effectiveness), RQ2 (Cost), RQ3 (Trade-offs), RQ4 (Explainability + Errors)
    \item \textbf{Experimental pipeline:} 5 bước từ data collection đến evaluation
    \item \textbf{Fairness controls:} Same hardware, seeds, hyperparameter protocol, evaluation metrics
    \item \textbf{Reproducibility protocol:} Code, data, environment, results package
    \item \textbf{Threats to validity:} Identified và mitigated
\end{enumerate}

Chương tiếp theo (Chapter 4) sẽ trình bày kết quả thực nghiệm, trả lời từng Research Question với data và visualizations.

% \include{chapters/04_results}
% \include{chapters/05_discussion}
% \include{chapters/06_conclusion}

% Bibliography
\newpage
\bibliographystyle{plain}
\bibliography{../literature/02_papers}
\addcontentsline{toc}{chapter}{Tài liệu tham khảo}

% Appendices
% \appendix
% \include{appendices/appendix_a}

\end{document}
